% =====================================================================
% Abstract (IEEEtran format)
% =====================================================================
\begin{abstract}
Robot motion planning routinely produces sequences of waypoints that
must be expanded into smooth time-parameterised trajectories before a
controller can execute them. This paper compares four numerical
interpolation methods drawn from Chapters~18 and~24 of Chapra and
Canale---linear interpolation, natural cubic splines, $C^4$ quintic
splines, and cubic B-spline approximation---under conditions
representative of two industrial robotics scenarios: a two-dimensional
mobile robot navigating an aisle and a two-degree-of-freedom planar
arm performing pick-and-place. Each method is implemented from the
textbook formulation rather than through library black-boxes,
validated against centered finite-difference derivatives
(Chapter~28), and evaluated on seven dataset variants that sweep over
waypoint density (5, 12, 20), geometric character (sharp 90$^{\circ}$
corners versus smooth sinusoid), and time-spacing (chord-length
versus uniform). The results challenge two pieces of conventional
wisdom: increasing waypoint density \emph{worsens} smoothness metrics
under chord-length parameterisation, and the natural-boundary quintic
spline produces \emph{higher} peak acceleration and root-mean-square
jerk than the cubic spline---an artefact entirely attributable to
the boundary condition rather than to the method, since cubic-derived
clamped boundary conditions reduce peak acceleration by 33\% and RMS
jerk by 64\%. The B-spline approximation produces the smoothest
derivatives at the cost of guaranteed waypoint deviation. We provide
a reusable Python implementation, an executable Jupyter notebook,
and a public dataset to support reproducibility.
\end{abstract}

\begin{IEEEkeywords}
B-splines, boundary conditions, cubic splines, finite differences,
interpolation, jerk minimisation, numerical methods, quintic splines,
robotics, trajectory planning, waypoint navigation.
\end{IEEEkeywords}
